\documentclass[11pt]{article}
\usepackage[a4paper,margin=28mm]{geometry}
\usepackage{amsmath,amssymb,amsthm,hyperref}
\newtheorem{proposition}{Proposition}
\title{A cubic-size barrier for the populated-cluster lifting method}
\author{Short asymptotic observation, proposed by the root agent}
\date{30 September 2026}
\begin{document}\maketitle
\begin{proposition}
Suppose two probability measures $\mu_A,\mu_B$ have disjoint ordered
clusters of support, each cluster belonging to only one measure.
Suppose each measure integrates polynomials of degree at most $m$
exactly as the uniform probability measure on $[0,1]$.
Then there are at least $m+1$ changes of color between consecutive
nonempty support clusters. In particular, each measure occupies at
least $\lfloor(m+2)/2\rfloor$ clusters.
\end{proposition}
\begin{proof}
The signed measure $\sigma=\mu_A-\mu_B$ annihilates every polynomial
of degree at most $m$. If its ordered nonempty support clusters had
$r\le m$ color changes, choose one point in each gap where a change
occurs. The product of the corresponding $r$ linear factors, multiplied
by a suitable global sign, has the sign of $\sigma$ on every support
cluster. Its integral against $\sigma$ is strictly positive, a
contradiction. There are therefore at least $m+2$ alternating color
blocks, requiring at least the stated number of clusters of each color.
\end{proof}

In a fully fair hybrid with these two variables and $m$ continuous
uniform variables, the polynomial exactness hypothesis follows simply
by restricting to either specified variable and the $m$ uniforms.
The method in \path{clustered_multi_exceptional_lifting.tex} requires
at least $m^2$ distinct nodes in each cluster. Consequently that
unmodified method necessarily uses at least
\[
 m^2\left\lfloor\frac{m+2}{2}\right\rfloor=\Omega(m^3)
\]
faces on each of the two specified dice. This is a restriction on that
particular lifting criterion. It does not give a cubic universal lower
bound and does not exclude two dice of order $m^2$ by another method.
\end{document}
